% Options for packages loaded elsewhere
\PassOptionsToPackage{unicode}{hyperref}
\PassOptionsToPackage{hyphens}{url}
\PassOptionsToPackage{dvipsnames,svgnames,x11names}{xcolor}
%
\documentclass[
  10pt,
  a4paper,
]{article}
\usepackage{amsmath,amssymb}
\usepackage{setspace}
\usepackage{iftex}
\ifPDFTeX
  \usepackage[T1]{fontenc}
  \usepackage[utf8]{inputenc}
  \usepackage{textcomp} % provide euro and other symbols
\else % if luatex or xetex
  \usepackage{unicode-math} % this also loads fontspec
  \defaultfontfeatures{Scale=MatchLowercase}
  \defaultfontfeatures[\rmfamily]{Ligatures=TeX,Scale=1}
\fi
\usepackage{lmodern}
\ifPDFTeX\else
  % xetex/luatex font selection
  \setmainfont[]{Liberation Serif}
  \setsansfont[]{Liberation Sans}
  \setmonofont[]{DejaVu Sans Mono}
\fi
% Use upquote if available, for straight quotes in verbatim environments
\IfFileExists{upquote.sty}{\usepackage{upquote}}{}
\IfFileExists{microtype.sty}{% use microtype if available
  \usepackage[]{microtype}
  \UseMicrotypeSet[protrusion]{basicmath} % disable protrusion for tt fonts
}{}
\makeatletter
\@ifundefined{KOMAClassName}{% if non-KOMA class
  \IfFileExists{parskip.sty}{%
    \usepackage{parskip}
  }{% else
    \setlength{\parindent}{0pt}
    \setlength{\parskip}{6pt plus 2pt minus 1pt}}
}{% if KOMA class
  \KOMAoptions{parskip=half}}
\makeatother
\usepackage{xcolor}
\usepackage[margin=22mm]{geometry}
\setlength{\emergencystretch}{3em} % prevent overfull lines
\providecommand{\tightlist}{%
  \setlength{\itemsep}{0pt}\setlength{\parskip}{0pt}}
\setcounter{secnumdepth}{-\maxdimen} % remove section numbering
\ifLuaTeX
\usepackage[bidi=basic]{babel}
\else
\usepackage[bidi=default]{babel}
\fi
\babelprovide[main,import]{korean}
\ifPDFTeX
\else
\babelfont{rm}[]{Liberation Serif}
\fi
% get rid of language-specific shorthands (see #6817):
\let\LanguageShortHands\languageshorthands
\def\languageshorthands#1{}
\usepackage{xeCJK}
\setCJKmainfont{Noto Serif CJK KR}
\setCJKsansfont{Noto Sans CJK KR}
\setCJKmonofont{Noto Sans Mono CJK KR}
\xeCJKsetup{CJKspace=true}
\usepackage{xurl}
\usepackage{fvextra}
\usepackage{fancyhdr}
\AtBeginDocument{\DefineVerbatimEnvironment{verbatim}{Verbatim}{breaklines=true,breakanywhere=true,fontsize=\small}}
\setlength{\emergencystretch}{4em}
\setlength{\headheight}{15pt}
\pagestyle{fancy}\fancyhf{}
\fancyhead[L]{\small 중성자별 구조와 비방사 진동}
\fancyfoot[C]{\thepage}
\renewcommand{\headrulewidth}{0.3pt}
\setcounter{tocdepth}{2}
\ifLuaTeX
  \usepackage{selnolig}  % disable illegal ligatures
\fi
\usepackage{bookmark}
\IfFileExists{xurl.sty}{\usepackage{xurl}}{} % add URL line breaks if available
\urlstyle{same}
\hypersetup{
  pdftitle={P2 · 숨은 binary와 외부 경로 찾기},
  pdflang={ko-KR},
  colorlinks=true,
  linkcolor={Maroon},
  filecolor={Maroon},
  citecolor={Blue},
  urlcolor={Blue},
  pdfcreator={LaTeX via pandoc}}

\title{P2 · 숨은 binary와 외부 경로 찾기}
\usepackage{etoolbox}
\makeatletter
\providecommand{\subtitle}[1]{% add subtitle to \maketitle
  \apptocmd{\@title}{\par {\large #1 \par}}{}{}
}
\makeatother
\subtitle{과제 큐시트}
\author{}
\date{}

\begin{document}
\maketitle

\setstretch{1.12}
\textbf{선행:} P1　 \textbf{독립 작업:} 0.5--1시간 (계획치)

\subsection{지난 실행의 흔적이 없어도 다시 만들 수
있는가?}\label{uxc9c0uxb09c-uxc2e4uxd589uxc758-uxd754uxc801uxc774-uxc5c6uxc5b4uxb3c4-uxb2e4uxc2dc-uxb9ccuxb4e4-uxc218-uxc788uxb294uxac00}

소스 폴더에 .o·.mod·실행파일이 섞여 있다. 그것들이 없어도 계산을 다시
만들 수 있어야 하는 이유를 검사 전에 설명하자.

먼저 한두 문장으로 답한 뒤 아래 설명을 읽자. 첫 예측이 틀렸거나 아직
모르겠다고 적는 것은 감점하지 않는다.

\subsection{필요한 관계와 읽을
위치}\label{uxd544uxc694uxd55c-uxad00uxacc4uxc640-uxc77duxc744-uxc704uxce58}

이 과제는 물리 방정식의 추가가 아니다. 동일한 소스와 공용 의존성을
보존하고, 다른 위치에서 실제로 빌드·실행할 수 있는지 확인한다.

\path|notes/chapters/ch11.md|, \path|tools/export_learner.py|,
\path|tools/verify_export.py|, \path|tools/make_calculation.py|

\subsection{계산으로
옮기기}\label{uxacc4uxc0b0uxc73cuxb85c-uxc62euxae30uxae30}

P1--P4의 자동 검사는 직접 구현한 learner의 독립 export 경로다. 제공
계산만 사용하는 독자는 별도 계산 폴더에서 같은 소스·의존성·실행 설명을
확인한다.

\begin{enumerate}
\def\labelenumi{\arabic{enumi}.}
\tightlist
\item
  learner export에서 binary·외부 경로·원본 소스 변경을 검사한다.
\item
  아래 명령은 SOURCE\_FILES.json을 가진 learner export에 사용한다.
  CALCULATION.json만 있는 폴더에는 적용하지 않는다.
\item
  거절된 항목의 실제 경로를 확인하고 과학 결과와 구별한다.
\end{enumerate}

\newpage

\subsection{실행해서 무엇을 볼
것인가}\label{uxc2e4uxd589uxd574uxc11c-uxbb34uxc5c7uxc744-uxbcfc-uxac83uxc778uxac00}

P0의 \texttt{export\_learner.py}로 만든 \texttt{../my\_solver}에
사용한다. 출력 파일은 새 이름으로 둔다.

\begin{verbatim}
python tools/verify_export.py --export ../my_solver --output ../source_check.json
\end{verbatim}

\subsection{결과를
설명하기}\label{uxacb0uxacfcuxb97c-uxc124uxba85uxd558uxae30}

숨은 binary나 원래의 절대 경로가 있으면 새 소스로 계산했다는 주장을
뒷받침하기 어렵다. 소스 보존 검사와 실행 성공은 구별한다.

\textbf{남길 증거.} 확인한 소스 목록, 외부 경로·binary 점검 결과와 그
한계

\textbf{판단 기준.} 2점: 실제 소스 보존·의존성·실행 기록을 연결하고,
계산하지 않은 물리 결과까지 검증했다고 말하지 않는다. 최초 예측의
정오보다, 결과를 본 뒤 이유를 고치고 다른 입력에도 적용하는지를 본다.

\subsection{막혔을 때}\label{uxb9c9uxd614uxc744-uxb54c}

\textbf{힌트 1.} 경로는 실행한 현재 폴더가 아니라 fpm.toml의 위치를
기준으로 해석한다.

\textbf{힌트 2.} SOURCE\_FILES.json에 기록된 파일과 실제 learner 파일의
내용을 비교한다.

\subsection{다음 만남에 다시
묻기}\label{uxb2e4uxc74c-uxb9ccuxb0a8uxc5d0-uxb2e4uxc2dc-uxbb3buxae30}

다음 회차에는 결과가 같은 두 폴더에서 실제 사용한 소스가 같은지 확인할
자료를 고른다. 파일명만 비교하지 않는다. 먼저 짧게 답한 뒤 본문과 자신의
기록을 열어 확인한다.

\textbf{범위.} 소스 보존과 build 성공은 별 전체의 물리 검증과 다르다.
작성하지 않은 함수는 새 폴더에서도 미완성이다.

\end{document}
